\documentclass[10.5pt]{article}
\usepackage{amsmath, amsfonts, amssymb,amsthm}
\usepackage[includeheadfoot]{geometry} % For page dimensions
\usepackage{fancyhdr}
\usepackage{enumerate} % For custom lists
\usepackage{tikz-cd}
\usepackage{graphicx}

\fancyhf{}
\lhead{Lie Algebras Notes}
\rhead{Tighe McAsey}
\pagestyle{fancy}

% Page dimensions
\geometry{a4paper, margin=1in}

\theoremstyle{definition}
\newtheorem{pb}{}
\newtheorem{definition}{Definition}
\newtheorem{proposition}{Proposition}
\newtheorem{example}{Example}
\newtheorem{theorem}{Theorem}
\newtheorem{remark}{Remark}
\numberwithin{example}{subsection}
\numberwithin{proposition}{subsection}
\numberwithin{pb}{subsection}
\usepackage{tikz-cd, stackengine}

% Commands:

\newcommand{\set}[1]{\left\{#1\right\}}
\newcommand{\gen}[1]{\langle#1\rangle}
\newcommand{\abs}[1]{\left\vert#1\right\vert}
\newcommand{\norm}[1]{\lvert\lvert#1\rvert\rvert}
\newcommand{\tand}{\text{ and }}
\newcommand{\tor}{\text{ or }}
\newcommand{\pd}{\frac{\partial}{\partial x_j}}
\newcommand{\px}{\frac{\partial}{\partial x}}
\newcommand{\py}{\frac{\partial}{\partial y}}
\newcommand{\pz}{\frac{\partial}{\partial z}}
\newcommand{\ppx}{\frac{\partial^2}{\partial x^2}}
\newcommand{\ppy}{\frac{\partial^2}{\partial y^2}}
\newcommand{\ppz}{\frac{\partial^2}{\partial z^2}}
\newcommand{\hess}{\operatorname{Hess}}
\newcommand{\inv}{\operatorname{inv}}
\newcommand{\spl}{\operatorname{\mathfrak{sl}}}
\newcommand{\spg}{\mathfrak{g}}
\newcommand{\sph}{\mathfrak{h}}
\newcommand{\spu}{\operatorname{\mathfrak{su}}}
\newcommand{\spo}{\operatorname{\mathfrak{so}}}


\begin{document}

    \section{Classification of Semi-Simple Lie Algebras}

    \subsection{Basic Examples}

    \subsubsection{\(\mathbf{A_\ell \sim \mathfrak{sl}_{\ell + 1}}\)}
    \begin{definition} \label{defsl}
        \begin{align}
            \mathfrak{sl}_{\ell + 1} = \mathfrak{gl}_{\ell + 1} \cap \set{\operatorname{tr} = 0}
        \end{align}
    \end{definition}

    The special linear lie algebra has dimension \((\ell + 1)^2 - 1\), an easy way to see this is to apply the submersion theorem to the determinant map on the lie group. A basis is given by \(\set{e_{i,j} \mid i \neq j} \cup \set{e_{ii} - e_{jj}}\).

    \begin{remark}
        A thought that comes to mind is that defining a lie algebra using structure constants (with some level of care about prime characteristic) the dimension is preserved, so in some sense the manifold structure for the real lie algebra can tell us about the dimensionality over an arbitrary field.
    \end{remark}

    \subsubsection{\(\mathbf{B_\ell} \sim \mathfrak{so}(2\ell + 1)\)}
        \begin{definition}\label{defbinformbl}
            Define the symmetric matrix \(s = \begin{pmatrix}
                1 & 0 & 0 \\ 0 & 0 & \mathbf{1_\ell} \\ 0 & \mathbf{1_\ell} & 0
            \end{pmatrix}\)
            which defines a symmetric bilinear form \(B\) on \(k^{2\ell + 1}\) defined as \(B(v,w) = vsw\).
        \end{definition}
        \begin{definition}\label{defsolodd}
            The na\"ive definition of \(\mathfrak{so}(2\ell + 1)\) is to take the tangent space to \(O(2\ell + 1)\), which gives us \(\mathfrak{gl} \cap \set{A^t = -A}\). This is equivalent to skew symmetric matrices with respect to the symmetric bilinear form defined by \(\mathbf{1_{2\ell + 1}}\). The drawback of this is it contains no diagonal elements, so that in particular, the Cartan subalgebra is not easily understood. Instead we define \(\mathfrak{so}(2\ell + 1)\) as skew symmetric matrices with respect to the binary form \(B\) given in \textbf{Defininition \ref{defbinformbl}}. Explicitly \(B(\xi(v),w) = -B(v,\xi(w))\).
        \end{definition}
        
        \begin{proposition}\label{equivsolodd}
            \begin{align}
                \mathfrak{so}(2\ell + 1) = \mathfrak{gl} \cap \set{\xi \mid s\xi = -\xi^t s} = \set{ \begin{pmatrix}
                0 & \mathbf{b_1} & \mathbf{b_2} \\ -\mathbf{b_1}^t & \mathbf{m} & \mathbf{n} \\ -\mathbf{b_2}^t & \mathbf{p} &-\mathbf{m}^t
                \end{pmatrix} \mid \mathbf{b_1},\mathbf{b_2},\mathbf{m},\mathbf{n},\mathbf{p}, -\mathbf{n}^t = \mathbf{n}, -\mathbf{p}^t = \mathbf{p}}
            \end{align}
        \end{proposition}
        \begin{proof}
            It suffices to show that \(\xi \in \mathfrak{so}(2\ell + 1) \iff s\xi = -\xi^t s\), since we read the third condition off of this one. If \(\xi \in \mathfrak{so}(2\ell + 1)\), then for all \(v,w \in k^{2\ell + 1}\)
            \begin{align}
                v^t\xi^tsw = \xi(v)^{t}sw = -v^ts\xi(w) = -v^ts\xi w
            \end{align}
            Since this holds for arbitrary \(v,w\) it is only possible if \(\xi^t s = -s \xi\). The converse is trivial.
        \end{proof}
        The third equivalent form in \textbf{Proposition \ref{equivsolodd}} lets us count the dimension of \(\mathfrak{so}(2\ell + 1)\) by elementary means of \(2\ell\) for the \(\mathbf{b_j}\), \(\ell^2\) for the \(\mathbf{m}\) and \(2\cdot\frac{\ell(\ell-1)}{2}\) for \(\mathbf{p},\mathbf{n}\) by counting the strictly upper triangular entries, in total the dimension is \(2\ell^2 + \ell\). This derivation gives us a basis.
    \subsubsection{\(\mathbf{C_\ell \sim \mathfrak{spl}(2\ell)}\)}

    \begin{definition}\label{defspeven}
        Let \(B\) be a skew symmetric non-degenerate bilinear form on the finite dimensional vectorspace \(V\), existence of such a form implies \(V\) has even dimension (\textbf{Proposition \ref{nondegenskewsym}}), at the risk of circular reasoning, the geometry of this is such a form defines an almost complex structure. Then we can take some isomorphism \(V \cong k^{2\ell}\)
        \begin{align}
            \mathfrak{spl}(2\ell) = \mathfrak{gl}(V)\cap\set{\xi \mid B(\xi(v),w) = -B(v,\xi(w))}
        \end{align}
        In practice we would like to define \(B\) by the skew symmetric matrix \(s = \begin{pmatrix} 0 & \mathbf{1_\ell} \\ -\mathbf{1_\ell} & 0 \end{pmatrix}\). Following the proof of \textbf{Proposition \ref{equivsolodd}} this is matrices \(\xi\) satisfying \(-\xi^t s = s \xi\).
    \end{definition}
    We can observe that
    \begin{align}
        \mathfrak{spl}(2\ell) = \set{\begin{pmatrix} \mathbf{m} & \mathbf{p} \\ \mathbf{q} & -\mathbf{m}^t \end{pmatrix}\mid \mathbf{p} = \mathbf{p}^t, \mathbf{q} = \mathbf{q}^t}
    \end{align}
    And the dimension is \(2 \ell^2 + \ell\)

    \subsubsection{\(\mathbf{D_\ell \sim \mathfrak{so}(2\ell)}\)}
    \begin{definition}\label{defsoleven}
        Once again, we don't take the naive skew symmetric approach, and define the lie algebra to be the maps satisfying \(B(\xi(v),w) = -B(v,\xi(w))\) where here \(B\) is defined by the matrix \(s = \begin{pmatrix} 0 & \mathbf{1_\ell} \\ \mathbf{1_\ell} & 0 \end{pmatrix}\). Thus we find that
        \begin{align}
            \mathfrak{so}(2\ell) = \set{\xi \in \mathfrak{gl}(2\ell) \mid -\xi^t s = s\xi} = \set{\begin{pmatrix} \mathbf{m} & \mathbf{p}\\\mathbf{q} & -\mathbf{m}^t \end{pmatrix} \mid \mathbf{p} = -\mathbf{p}^t,\mathbf{q} = \mathbf{q}^t}
        \end{align}
    \end{definition}
    Once again it is easy to see the dimension is \(2\ell^2 + \ell\) by using the explicit form and counting skew symmetric matrices.

    
    \subsection{Definitions and Major Theorems}

    \begin{definition}
        An ideal of a Lie algebra is \textcolor{red}{TODO}
    \end{definition}

    \begin{definition}
        The center \(\mathfrak{z}\) of a Lie algebra is \textcolor{red}{TODO}
    \end{definition}

    \begin{definition}
        A subalgebra of a Lie algebra is \textcolor{red}{TODO}
    \end{definition}

    \begin{definition}
        A Lie algebra homomorphism is \textcolor{red}{TODO}
    \end{definition}

    \begin{definition}
        A representation of a Lie algebra is \textcolor{red}{TODO}
    \end{definition}

    \begin{definition}
        The derived subalgebra of a Lie algebra is \textcolor{red}{TODO}
    \end{definition}

    \begin{definition}
        The normalizer \(N_\mathfrak{g}(\mathfrak{h})\) of a subalgebra \(\mathfrak{h}\) of a Lie algebra \(\mathfrak{g}\) is \textcolor{red}{TODO}
    \end{definition}

    \begin{definition}\label{Derivation}
        A derivation \(D: \mathfrak{g} \to \mathfrak{g}\) is an endomorphism satisfying 
        \begin{equation}
            D[\xi,\eta] = [D(\xi),\eta] + [\xi,D(\eta)]
        \end{equation}
    \end{definition}

    \begin{definition}\label{expofder}
        In characteristic zero, the exponential of a nilpotent derivation (i.e. \(D^j = 0\) for some \(j\)) is the formal power series \(\exp(D) = \sum_0^\infty D^j/j!\). By the Liebniz rule \(\exp(D)\) is a derivation.
    \end{definition}

    \begin{definition}\label{defkilling} The killing form \(\kappa\) is a bilinear form on \(\mathfrak{g}\) defined as
            \begin{align}
                \kappa: \mathfrak{g} \times \mathfrak{g} &\to \mathfrak{g} \\
                (\xi,\eta) &\mapsto \operatorname{tr}(\operatorname{ad}_\xi\circ\operatorname{ad}_\eta)
            \end{align}
        \end{definition}
        Bilinearity follows from linearity of the arguments in the adjoint representation (i.e. bilinearity of lie bracket).

        \begin{definition}\label{defadinvariant}
            A bilinear form \(B(-,-)\) on \(\mathfrak{g}\) is called \(\operatorname{ad}\) invariant, or simply invariant, when
            \begin{align}
                B([x,\xi],\eta) + B(\xi,[x,\eta]) = 0
            \end{align}
            for all \(x,\xi,\eta\).
        \end{definition}

        \begin{definition}\label{definng}
            The Lie Algebra of inner derivations of \(\mathfrak{g}\), denoted \(\operatorname{Inn}(\mathfrak{g})\) is given by
            \begin{align}
                \operatorname{Inn}(\mathfrak{g}) = \set{\operatorname{ad}_\xi \mid \xi \in \mathfrak{g}}
            \end{align}
            The Lie bracket is inhereted from \(\operatorname{Der}(\mathfrak{g})\), which is inhereted from \(\mathfrak{gl}(\mathfrak{g})\)
        \end{definition}
    
    \subsection{The Killing Form}

        \textbf{See Definition \ref{defkilling}.}

        \begin{proposition}\label{killingsymmetric}
            The killing form is a symmetric bilinear form.
        \end{proposition}
        \begin{proof}
            \(\operatorname{tr}\) is symmetric.
        \end{proof}

        \begin{proposition}\label{killingautinvariant}
            \(\kappa\) is automorphism invariant. I.e. Let \(\phi \in \operatorname{Aut}(\mathfrak{g})\), then
            \begin{align}
                \kappa(\phi(\xi),\phi(\eta)) = \kappa(\xi,\eta)
            \end{align}
        \end{proposition}
        \begin{proof}
            \begin{align}
                \operatorname{ad}_{\phi(\xi)} = \phi\circ \operatorname{ad}_\xi \circ \phi^{-1}
            \end{align}
            and trace is basis independent.
        \end{proof}

        \begin{proposition}\label{killingadinvariant}
            The killing form is an invariant binary form \textbf{See Definition \ref{defadinvariant}}. More generally it is invariant (in the same sense) under all derivations.

            \begin{proof}
                Choose a path in \(\operatorname{Aut}(\mathfrak{g})\) which differentiates to the desired inner derivation. Then differentiate equation for automorphism invariance using proposition \ref{dergautg}, apply bilinearity of the killing form to take the derivative.
            \end{proof}

            A second proof is much less elegant, put proves the result for general derivations.
            \begin{proof}
                By the definition of derivations,
                \begin{align}\label{adder}
                    \operatorname{ad}_{D(x)} = D\circ \operatorname{ad}_x  - \operatorname{ad}_x\circ D
                \end{align}
                Thus we find that
                \begin{align}
                    \kappa(D(\xi),\eta) &= \operatorname{tr}(\operatorname{ad}_{D(\xi)}\circ\operatorname{ad}_\eta) \\
                    & \overset{\eqref{adder}}{=} \operatorname{tr}(D\circ \operatorname{ad}_\xi\circ \operatorname{ad}_\eta)  - \operatorname{tr}(\operatorname{ad}_\xi\circ D\circ \operatorname{ad}_\eta) \\
                    &= - \left(\operatorname{tr}(\operatorname{ad}_\xi\circ D \circ \eta) - \operatorname{tr}(\operatorname{ad}_\xi \circ \operatorname{ad}_\eta \circ D) \right) \label{justify1} \\
                    &=\operatorname{tr}(\operatorname{ad}_\xi \circ(D\circ \operatorname{ad}_\eta  - \operatorname{ad}_\eta\circ D)) \\
                    & \overset{\eqref{adder}}{=} - \operatorname{tr}(\operatorname{ad}_\xi \circ \operatorname{ad}_{D(\eta)}) = -\kappa(\xi,D(\eta))
                \end{align}
                Where \eqref{justify1} follows from cyclicity of trace.
            \end{proof}
        \end{proposition}

        \begin{proposition}\label{derivationsareinner}
            If the killing form \(\kappa\) of \(\mathfrak{g}\) is nondegenerate, then every derivation of \(\mathfrak{g}\) is inner. I.e.
            \begin{align}
                \operatorname{Der}(\mathfrak{g}) = \operatorname{Inn}(\mathfrak{g})
            \end{align}

            \begin{proof}
                Since \(\kappa\) is nondegenerate, the map \(\mathfrak{g} \to \mathfrak{g}^*\) via \(\xi \mapsto \kappa(\xi,-)\) is an isomorphism. Now let \(D\) be a derivation, we define the functional
                \begin{align}
                    \eta \mapsto \operatorname{tr}(D\circ \operatorname{ad}_\eta)
                \end{align}
                which must be equal to \(\eta \mapsto \kappa(\xi,\eta)\) for some \(\xi\). Fixing \(x \in \mathfrak{g}\) we find that
                \begin{align}
                    \eta \mapsto \kappa((D - \operatorname{ad}_\xi)(x),\eta) &= \kappa(D(x),\eta) - \kappa(\operatorname{ad}_\xi(x),\eta) \\
                    &= \operatorname{tr}(\operatorname{ad}_{D(x)}\circ \operatorname{ad}_\eta) - \operatorname{tr}(\operatorname{ad}_{\operatorname{ad}_{\xi}(x)}\circ \operatorname{ad}_\eta) \\
                    &\overset{\eqref{adder}}{=} \operatorname{tr}([D,\operatorname{ad}_x]\circ \operatorname{ad}_\eta) - \operatorname{tr}([\operatorname{ad}_\xi,\operatorname{ad}_x]\circ \operatorname{ad}_\eta) \\
                    &= \operatorname{tr}(D\circ[\operatorname{ad}_x,\operatorname{ad}_\eta]) - \operatorname{tr}(\operatorname{ad}_\xi\circ[\operatorname{ad}_x,\operatorname{ad}_\eta]) \label{justify2} \\
                    &\overset{\eqref{adder}}{=} \operatorname{tr}(D\circ \operatorname{ad}_{\operatorname{ad}_x(\eta)}) - \kappa(\xi,\operatorname{ad}_x(\eta)) \\
                    &= \operatorname{tr}(0) = 0
                \end{align}
                Where \eqref{justify2} follows from cyclicity of trace.
            \end{proof}
        \end{proposition}

    \section{Lie Groups}
    \subsection{General Theory}
        \begin{theorem}\label{thmclosedsubgroup}
            A closed subgroup of a Lie Group is a Lie Group
        \end{theorem}
    \subsection{Lie Groups For studying Lie Algebras}
        \begin{proposition}\label{dergautg}
                \(\operatorname{Aut}(\mathfrak{g})\) is a Lie group, with Lie algebra \(\operatorname{Inn}(\mathfrak{g})\).
                \begin{proof}
                    \(\operatorname{GL}(\mathfrak{g})\) is certainly a Lie group, and \(\operatorname{Aut}(\mathfrak{g})\) is a subgroup which is closed by continuity of \([-,-]\) so \textbf{Theorem \ref{thmclosedsubgroup}} applies. By bilinearity, differentiation of the equation if \(\phi(t)\) is a path in \(\operatorname{Aut}(\mathfrak{g})\) with \(\phi(0) = \operatorname{id}\) then
                    \begin{align}
                        \left.\frac{d}{dt}\right\vert_{t=0}\phi(t)([\xi,\eta]) &= \left.\frac{d}{dt}\right\vert_{t=0}[\phi(t)(\xi),\phi(t)(\eta)] \\
                        &= \left[\left.\frac{d}{dt}\right\vert_{t=0}\phi(t)(\xi),\eta\right] + \left[\xi,\left.\frac{d}{dt}\right\vert_{t=0}\phi(t)(\eta)\right]
                    \end{align}
                \end{proof}
        \end{proposition}
        The converse direction follows since the automorphism equation integrates the derivation equation.

        \section{Exercises}
        \subsection{Humphreys Ch1.}
        \begin{pb}
            Compute the structure constants of the cross product on \(\mathbb{R}^3\)
            \begin{proof}
                \(a_{12}^3 = 1, a_{31}^2 = 1, a_{23}^1 = 1\), else zero.
            \end{proof}
        \end{pb}
        \begin{pb}
            Verify that the bilinearity and "antisymmetry" axioms alongside
            \begin{align}
                [x,y] = z, [x,z] = y, [y,z] = 0
            \end{align}
            define a Lie algebra.
            \begin{proof}
                It suffices to check the Jacobi equations
                \begin{align}
                    [x,[y,z]] - [y,[x,z]] + [z,[x,y]] = 0 - [y,y] + [z,z] = 0
                \end{align}
            \end{proof}
        \end{pb}
        \begin{pb}
            Let \(e,h,f = e_{12}, e_{11} - e_{22}, e_{21}\) be the standard basis of \(\mathfrak{sl}(2)\), compute the matrices \(\operatorname{ad}_{-}\)
            \begin{proof}
                \begin{align}
                    \operatorname{ad}_e &= -2e_{12} + e_{23} \\
                    \operatorname{ad}_f &= -e_{21} + 2e_{32} \\
                    \operatorname{ad}_h &= 2e_{11} - 2e_{33}
                \end{align}
            \end{proof}
        \end{pb}
        \begin{pb}
            Find a linear lie algebra \(\mathfrak{g}\) with \(\mathfrak{g} = \operatorname{span}_k\set{x,y}\) and \([x,y] = x\).
            \begin{proof}
                We exploit the symmetry of the adjoint representation \(\operatorname{ad}_{[x,y]} = [\operatorname{ad}_x,\operatorname{ad}_y]\), and take \(x,y\)to their adjoint representations. This works since \(\operatorname{Inn}(\mathfrak{g})\) is also a 2 dimensional lie algebra.
                \begin{align}
                    x = e_{12}, y = e_{11}
                \end{align}
            \end{proof}
        \end{pb}
        \begin{pb}
            Compute the dimensions of the diagonal, upper triangular and strictly upper triangular lie algebras. Verify that \([\mathfrak{d},\mathfrak{n}] = \mathfrak{n}\).
            \begin{proof}
                The dimensiona are \(n, \frac{n(n+1)}{2}, \frac{n(n-1)}{2}\). The bracket condition follows from \([e_{jj},e_{k,r}] = \delta_{jk}e_{jr} - \delta_P{rj}e_{kj}\) But by assuming strictly upper triangular \(j = k\) implies \(r > j\) and \(r = j\) implies \(j > k\).
            \end{proof}
        \end{pb}
        \begin{pb}
            Suppose \(x \in \mathfrak{gl}(n,k)\) has \(n\) distinct eigenvalues \(a_j\). Show that the eigenvalues of \(\operatorname{ad}_x\) are \(a_i - a_j\).

            \begin{proof}
                Work in the basis of eigenvectors for \(x\), then
                \begin{align}
                    \operatorname{ad}_x(e_{ij})(e_k) = \begin{cases}
                        0 & k \neq j \\
                        x(e_i) - e_{ij}(a_je_j) = (a_i - a_j)e_i & k = j
                    \end{cases}
                \end{align}
                Since the \(e_{ij}\) form a basis, each is an eigenvector with eigenvalue \(a_i - a_j\).
            \end{proof}
        \end{pb}
        \begin{pb}
            Show that \(\mathfrak{gl}(n,k) = \mathfrak{sl}(n,k) + \operatorname{span}_k\set{\mathbf{1_n}}\) where char \(k\) is not a prime dividing \(n\)

            \begin{proof}
                Let \(\xi \in \mathfrak{gl}(n)\), then \(\xi - n^{-1}\operatorname{tr}(\xi)\mathbf{1_n} \in \mathfrak{sl}(n)\)
            \end{proof}
        \end{pb}
        \begin{pb}
            Verify the dimension of \(\mathbf{D_\ell}\) is \(\ell^2 + \ell\)
            \begin{proof}
                See definition \ref{defsoleven}.
            \end{proof}
        \end{pb}
        \begin{pb}
            Show that in characteristic zero, the classical lie algebras are equal to their derived subalgebras.

            \begin{proof}
                Notice that by defining the orthogonal algebras as we did, all of the classical lie algebras have diagonal elements. Thus \(\operatorname{tr}: \mathfrak{g} \times \mathfrak{g} \to k\) by \(\xi,\eta \mapsto \operatorname{tr}(\xi \eta)\) is a nondegenerate bilinear form. Moreover by cyclicity it is an invariant bilinear form. This tells us that
                \begin{align}
                    \operatorname{tr}([\mathfrak{g},\mathfrak{g}],-) = \operatorname{tr}(\mathfrak{g},[\mathfrak{g},-])
                \end{align}
                Thus \(\xi \in [\mathfrak{g},\mathfrak{g}]^\perp \implies \xi \in \mathfrak{z}(\mathfrak{g})\). Semisimple lie algebras have trivial center so we are done.
            \end{proof}
        \end{pb}
        \begin{pb}
            Exhibit the exceptional isomorphisms \(A_1 \cong B_1 \cong C_1\), show \(D_1\) has dimension 1, show \(B_2 \cong C_2\), \(D_3 \cong A_3\). What about \(D_2\)?

            \begin{proof}
                \(\mathbf{A_1 \cong B_1.}\) Consider \(\operatorname{ad}: \mathfrak{sl}(2) \to \mathfrak{gl}(3)\), on \(\mathfrak{sl}(2)\), \(\operatorname{tr}\) is a nondegenerate symmetric bilinear form, moreover, cyclicity gives us \(\operatorname{tr}(\operatorname{ad}_\xi(x),y) = -\operatorname{tr}(x,\operatorname{ad}_\xi(y))\). Thus by equivalence of bilinear forms (we are using here the field is algebraically closed), we get \(\operatorname{ad}(\mathfrak{sl}(2)) \subset \mathfrak{so}(3)\). Note that this proof fails due to the killing form on \(\mathfrak{sl}(2,\mathbb{R})\) not being positive definite.

                \(\mathbf{A_1 \cong C_1.}\) \(-\xi^ts = s\xi\) if and only if \(\operatorname{tr}(\xi) = 0\)

                \(\mathbf{D_1 \cong U(1)}\) By direct computation \(D_1 = \set{\begin{pmatrix} t & 0 \\ 0 & t^{-1} \end{pmatrix}}\)

                \textcolor{red}{TODO: Prove the remaining exceptional isomorphisms}
            \end{proof}
        \end{pb}
        \begin{pb}
            Verify the commutator of two derivations is a derivation, show that the composition of derivations on an algebra need not be one.
            
            \begin{proof}
                \begin{align}
                    [D,E][\xi,\eta] &= D([E\xi,\eta] + [\xi,E\eta]) - E([D\xi,\eta] + [\xi,D\eta]) \\
                    &= [DE\xi,\eta] + [E\xi,D\eta] + [D\xi,E\eta] + [\xi,DE\eta] - ([E\xi,D\eta] + [D\xi,E\eta]) - [ED\xi,\eta] - [\xi,ED\eta] \\
                    &= [[D,E]\xi,\eta] + [\xi,[D,E]\eta]
                \end{align}

                Take the example of the formal derivative on \(k[X]\).
            \end{proof}
        \end{pb}
        \begin{pb}
            Show that \(\set{\eta \in \mathfrak{g} \mid \exists \lambda \in k, \operatorname{ad}_\xi(\eta) = \lambda \eta} \subset \mathfrak{g}\) is a subalgebra.

            \begin{proof}
                Apply the Jacobi identity to get
                \begin{align}
                    \operatorname{ad}_\xi[\eta,x] &= [\eta,\operatorname{ad}_\xi x] + [x, -\operatorname{ad}_\xi \eta] \\
                    &=[\eta, \lambda_x x] + [x,-\lambda_\eta \eta] \\
                    &= (\lambda_x + \lambda_\eta)[\eta,x]
                \end{align}
            \end{proof}
        \end{pb}

        \subsection{Humphreys Ch2.}

        \begin{pb}
            Show that \(\operatorname{Inn}(\mathfrak{g}) \subset \operatorname{Der}(\mathfrak{g})\) is an ideal

            \begin{proof}
                \textcolor{red}{TODO}
            \end{proof}
        \end{pb}
        \begin{pb}
            Show that \(\mathfrak{sl}(n) = [\mathfrak{gl}(n),\mathfrak{gl}(n)]\)

            \begin{proof}
                \textcolor{red}{TODO}
            \end{proof}
        \end{pb}
        \begin{pb}
            Show the center of \(\mathfrak{gl}(n)\) is diagonal matrices, show the center of \(\mathfrak{sl}(n)\) is trivial, but in characteristic \(p\vert n\) is diagonal matrices.

            \begin{proof}
                \textcolor{red}{TODO}
            \end{proof}
        \end{pb}
        \begin{pb}
            Show there is a unique isomorphism class of lie algebra \(\mathfrak{g}\) with dimension 3 and \(\mathfrak{g}' = [\mathfrak{g},\mathfrak{g}] \subset \mathfrak{z}(\mathfrak{g})\), and \(\dim \mathfrak{g}' = 1\).

            \begin{proof}
                \textcolor{red}{TODO}
            \end{proof}
        \end{pb}
        \begin{pb}
            Show that \(\dim \mathfrak{g} = 3\) and \(\mathfrak{g}' = \mathfrak{g}\) implies \(\mathfrak{g}\) is simple.

            \begin{proof}
                \textcolor{red}{TODO}
            \end{proof}
        \end{pb}
        \begin{pb}
            Show \(\mathfrak{sl}(3)\) is simple unless characteristic \(k = 3\).
        \end{pb}
        \begin{pb}
            Show that strictly upper triangular and diagonal matrices are self normalizing in \(\mathfrak{gl}\), show that upper triangular matrices are normalized by strictly upper triangular matrices.

            \begin{proof}
                \textcolor{red}{TODO}
            \end{proof}
        \end{pb}
        \begin{pb}
            If \(k\) is a field of characteristic zero then the diagonal matrices are self normalizing subsets of the classical lie algebras.

            \begin{proof}
                \textcolor{red}{TODO}
            \end{proof}
        \end{pb}
        \begin{pb}
            \textcolor{red}{TODO}
        \end{pb}
        \begin{pb}
            for \(g \in GL(n)\) consider \(\mathfrak{sl}(n) \to \mathfrak{sl}(n)\) via \(x \mapsto -gx^tg^{-1}\), show this is an automorphism of \(\mathfrak{sl}(n)\). When \(n=2\) and \(g = 1\) show this is an inner automorphism.

            \begin{proof}
                \textcolor{red}{TODO}
            \end{proof}
        \end{pb}
        \begin{pb}
            Prove that for \(g \in O(n)\) the map \(x \mapsto gxg^{-1}\) defines an automorphism of an orthogonal lie algebra.
        \end{pb}


        \section{Appendix}
        \begin{proposition}\label{bilinearformequiv}
            All non-degenerate skew symmetric Bilinear forms are equivalent.

            Over an algebraically closed field, all non-degenerate symmetric binary forms are equivalent. The obstruction is existence of nonsquares (i.e. over \(\mathbb{R}\) we classify which signs we act positively or negatively, over \(\mathbb{Q}\) there are many ways forms can be non equivalent).
        \end{proposition}
        \begin{proposition}\label{nondegenskewsym}
            A nondegenerate skew symmetric bilinear form can only exist in even dimensions.

            \begin{proof}
                Assume that \(\dim V = 2\ell + 1\) has a nondegenerate bilinear form \(B\). Then by induction and graham schmidt, we get a basis for a \(2\ell\) dimensional subspace \(e_1,\hdots,e_\ell,d_1,\hdots,d_\ell\) with \(B(e_j,d_j) = 1\) and all other products zero. Then there is another linearly independent element \(f\), define \(f'\) as follows
                \begin{align}
                    f' = f + \sum_j B(f,e_j)d_j - B(f,d_j)e_j
                \end{align}
                Then \(B(f',-) = 0\) so that \(f' = 0\) contradicting \(V\) being \(2 \ell + 1\) dimensional. Construction of the basis from induction follows a similar technique so I won't make it explicit.
            \end{proof}
        \end{proposition}
        
\end{document}